\documentclass[11pt, a4paper]{article}
\usepackage[a4paper, top=2.5cm, bottom=2.5cm, left=2cm, right=2cm]{geometry}
\usepackage{amsmath,amssymb,amsthm,microtype}
\usepackage{hyperref}

\theoremstyle{definition}
\newtheorem{problem}{Problem}
\theoremstyle{plain}
\newtheorem{theorem}{Theorem}
\newtheorem{lemma}{Lemma}

\begin{document}

\title{\vspace{-1.5cm}\textbf{Regional Quantitative Problem-Solving Circle} \\ \Large Solution Key 6: Stochastic Modeling, Markov Chains \& Synthesis (Weeks 11--12)}
\author{\textbf{Lead Architect:} Mohamed Jad Srifi}
\date{\textbf{Term:} Fall 2025 -- Spring 2026}
\maketitle

\section*{Rigorous Proofs \& Structural Solutions}

\begin{problem}[Gambler's Ruin / First-Step Difference Equations --- Putnam Caliber]
Derive the boundary value recurrence for ruin probability $R_k$, analyze telescoping differences $(R_{k+1} - R_k)$, and derive exact closed-form expressions for $p = 1/2$ and $p \ne 1/2$.
\end{problem}

\begin{proof}
Let $R_k$ be the probability of ruin starting from state $k$. By First-Step Analysis, the process transitions to state $k+1$ with probability $p$ and to state $k-1$ with probability $q$. The recurrence is:
\[
R_k = p R_{k+1} + q R_{k-1}
\]
The absorbing boundary conditions are $R_0 = 1$ (immediate ruin) and $R_N = 0$ (immediate victory).
Because $p + q = 1$, we can rewrite the left side as $(p + q)R_k$:
\begin{align*}
p R_k + q R_k &= p R_{k+1} + q R_{k-1} \\
p(R_{k+1} - R_k) &= q(R_k - R_{k-1}) \\
R_{k+1} - R_k &= \frac{q}{p} (R_k - R_{k-1})
\end{align*}
Define the first-order difference $\Delta_k = R_{k+1} - R_k$. The relation establishes a geometric progression: $\Delta_k = (\frac{q}{p}) \Delta_{k-1}$.
By recursive substitution down to the base difference $\Delta_0 = R_1 - R_0$:
\[
\Delta_i = \left(\frac{q}{p}\right)^i (R_1 - R_0)
\]
We isolate $R_k$ by summing the telescoping differences from index $0$ to $k-1$:
\begin{align*}
\sum_{i=0}^{k-1} (R_{i+1} - R_i) &= R_k - R_0 \\
R_k - 1 &= \sum_{i=0}^{k-1} \left(\frac{q}{p}\right)^i (R_1 - 1)
\end{align*}
\textbf{Case 1 ($p = 1/2 \implies q/p = 1$):} 
The summation reduces to summing the constant $1$ exactly $k$ times.
\[
R_k - 1 = k(R_1 - 1)
\]
We apply the upper boundary $R_N = 0$ to solve for $R_1$: $0 - 1 = N(R_1 - 1) \implies (R_1 - 1) = -1/N$.
Substituting back: $R_k = 1 - \frac{k}{N}$.

\textbf{Case 2 ($p \ne 1/2 \implies q/p \ne 1$):}
The summation is a standard finite geometric series:
\[
R_k - 1 = (R_1 - 1) \frac{1 - (q/p)^k}{1 - (q/p)}
\]
We apply the upper boundary $R_N = 0$ to isolate $(R_1 - 1)$:
\[
-1 = (R_1 - 1) \frac{1 - (q/p)^N}{1 - (q/p)} \implies (R_1 - 1) = - \frac{1 - (q/p)}{1 - (q/p)^N}
\]
Substituting back into the $R_k$ equation yields the closed-form ruin probability:
\[
R_k = 1 - \frac{1 - (q/p)^k}{1 - (q/p)^N} = \frac{(q/p)^k - (q/p)^N}{1 - (q/p)^N}
\]
\end{proof}

\vspace{0.5cm}

\begin{problem}[Expected Hitting Times on Symmetric Cyclic Graphs --- MIT 6.042 Standard]
Derive the recurrence $E_k = 1 + 0.5 E_{k-1} + 0.5 E_{k+1}$, evaluate $\Delta^2 E_k = -2$, and prove $E_k = k(N - k)$.
\end{problem}

\begin{proof}
Let $E_k$ be the expected time to reach the absorbing node $0$ starting from node $k$. By First-Step Analysis, taking one step consumes $1$ time unit, after which the particle transitions to $k-1$ or $k+1$ with equal probability $1/2$. Thus:
\[
E_k = 1 + \frac{1}{2}E_{k-1} + \frac{1}{2}E_{k+1} \quad \text{for } 1 \le k \le N-1
\]
The boundary conditions are $E_0 = 0$ and $E_N = 0$ (since node $N$ is topologically identical to node $0$ on the cycle).
We rearrange the recurrence to isolate the discrete differences. Multiply by $2$:
\begin{align*}
2E_k &= 2 + E_{k-1} + E_{k+1} \\
E_k - E_{k-1} - (E_{k+1} - E_k) &= 2
\end{align*}
Define the first-order forward difference $\Delta E_k = E_{k+1} - E_k$. Substituting this yields:
\begin{align*}
- (\Delta E_k - \Delta E_{k-1}) &= 2 \\
\Delta^2 E_k &= -2
\end{align*}
The second-order finite difference is a strictly negative constant ($-2$). In discrete calculus, a constant second difference implies the solution sequence $E_k$ must be a purely quadratic polynomial of the form:
\[
E_k = Ak^2 + Bk + C
\]
We apply our boundary conditions to solve for the coefficients. 
From $E_0 = 0$, we immediately derive $C = 0$.
To satisfy $\Delta^2 E_k = -2$, the leading coefficient of the quadratic must be half of the second difference (analogous to the continuous second derivative). Thus, $A = -2 / 2 = -1$. The equation simplifies to $E_k = -k^2 + Bk$.
From the topological wrap-around boundary $E_N = 0$:
\[
-N^2 + BN = 0 \implies B = N
\]
Substituting $B$ into our polynomial provides the exact expected hitting time:
\[
E_k = -k^2 + Nk = k(N - k)
\]
\end{proof}

\vspace{0.5cm}

\begin{problem}[Synthesis: Coupling \& Stationary Distribution Convergence]
Formulate the dual-particle system as a single absorbing Markov chain and derive the exact expected meeting time.
\end{problem}

\begin{proof}
The total state space of the two particles is strictly vast, but due to the complete symmetry of the complete graph $K_M$, we can collapse the state space into a binary Markov chain based strictly on spatial intersection. Let the states be $S_0$ (particles occupy distinct vertices) and $S_1$ (particles occupy the same vertex). State $S_1$ is absorbing.
We must calculate $p$, the probability of transitioning from $S_0$ to $S_1$ in a single discrete step. Let the particles originate at vertices $A$ and $B$. There are exactly three mutually exclusive pathways for a collision to occur:
\begin{enumerate}
    \item Particle A stays at $A$, and Particle B moves to $A$. The joint probability is:
    \[ \frac{1}{2} \times \frac{1}{2(M-1)} = \frac{1}{4(M-1)} \]
    \item Particle B stays at $B$, and Particle A moves to $B$. Symmetrically, the probability is:
    \[ \frac{1}{2(M-1)} \times \frac{1}{2} = \frac{1}{4(M-1)} \]
    \item Both Particle A and Particle B move simultaneously to a third, identical vertex $v \notin \{A, B\}$. There are exactly $(M - 2)$ such available vertices. The joint probability is:
    \[ (M - 2) \times \left( \frac{1}{2(M-1)} \right) \times \left( \frac{1}{2(M-1)} \right) = \frac{M-2}{4(M-1)^2} \]
\end{enumerate}
Since these pathways are mutually exclusive, we sum them to find the total transition probability $p$:
\begin{align*}
p &= \frac{1}{4(M-1)} + \frac{1}{4(M-1)} + \frac{M-2}{4(M-1)^2} \\
p &= \frac{(M-1) + (M-1) + (M-2)}{4(M-1)^2} \\
p &= \frac{3M - 4}{4(M-1)^2}
\end{align*}
Because the system is memoryless and the collision probability $p$ is strictly constant at every step in $S_0$, the number of steps $T$ required to hit the absorbing state $S_1$ perfectly follows a geometric distribution.
The expected value of a geometric random variable is the reciprocal of its success probability:
\[
\mathbb{E}[T] = \frac{1}{p} = \frac{4(M-1)^2}{3M - 4}
\]
Thus, the expected time for the independent random walks to intersect is strictly bounded.
\end{proof}

\end{document}
